\chapter{Masuk dan Mengenal Pandora}

Bab ini membahas langkah pertama yang Anda lakukan setiap kali memakai Pandora. Anda akan masuk dengan akun UGM, membaca status layanan di bagian atas layar, dan mengenali kelompok panel yang akan Anda pakai di bab-bab berikutnya. Rina memulai justru dari sini, sebab ia tidak ingin menunggu di tengah rapat tanpa tahu apakah layanan siap dipakai.

\begin{indikator}
Setelah mempelajari bab ini, peserta mampu:
\begin{enumerate}
  \item masuk ke Pandora dengan akun UGM yang sah;
  \item membaca status layanan di bagian atas layar dan menentukan apakah fitur AI siap dipakai;
  \item menyebutkan fungsi setiap kelompok panel pada layar utama.
\end{enumerate}
\end{indikator}

\section{Masuk dengan Akun UGM}
Pandora menggunakan login Google. Akun yang diterima adalah akun dengan domain \texttt{ugm.ac.id} atau \texttt{mail.ugm.ac.id}. Setelah memilih akun, Anda diarahkan kembali ke halaman utama Pandora dan namanya tampil di layar.

Akun Anda menentukan kepemilikan hasil. Setiap transkrip, notulensi, dan dokumen yang Anda buat hanya tampil di riwayat akun itu. Karena itu, jangan membagikan komputer yang sedang terbuka di halaman Pandora kepada orang lain, dan keluar dari akun ketika selesai bekerja di komputer bersama.

\section{Membaca Status Layanan}
Bagian atas layar memperlihatkan status komponen yang dibutuhkan. Status ini membantu Anda tahu apa yang bisa dipakai sebelum memulai pekerjaan. Tabel~\ref{tab:status} merangkum arti setiap status.

\begin{table}[H]
\caption{Status layanan pada bagian atas layar}
\label{tab:status}
{\small
\begin{tabularx}{\linewidth}{@{}P{4.2cm}Y@{}}
\toprule
\rowcolor{ugmpanel}\thd{Komponen} & \thd{Arti dan tindakan} \\
\midrule
Pengolah audio & Komponen yang mengubah format rekaman. Bila bermasalah, transkripsi tidak bisa dimulai. Laporkan kepada pengelola. \\
Model transkripsi & Model pengenal suara tersedia. Bila tidak tersedia, pilih model lain yang tersedia atau minta pengelola memeriksanya. \\
Layanan AI & Diperlukan untuk merapikan transkrip, notulensi, dan dokumen resmi. Bila belum siap, tombol AI tidak dapat diklik. Transkripsi tetap bisa dipakai. \\
\bottomrule
\end{tabularx}
}
\end{table}

Tombol untuk merapikan transkrip, membuat notulensi, dan menyusun dokumen resmi hanya aktif bila layanan AI siap dan sebuah folder hasil sudah dipilih. Sebelum mengandalkan fitur AI untuk rapat yang penting, periksa status ini terlebih dahulu. Bila status tidak sesuai, catat waktunya dan beri tahu pengelola sebelum rapat berjalan.

\section{Mengenal Tata Letak Layar}
Layar utama Pandora terbagi menjadi beberapa kelompok panel. Gambar~\ref{fig:tata-letak} menunjukkan posisi kelompok utama, dan Tabel~\ref{tab:panel} menjelaskan fungsinya.

\begin{figure}[H]
\centering
\begin{tikzpicture}
  \node[kotakgelap, minimum width=15cm, minimum height=9mm] (h) at (0,0) {Bagian atas: status layanan};
  \node[kotak, minimum width=7.3cm, minimum height=24mm] (u) at (-3.9,-1.7) {Panel unggah dan pengaturan\\(berkas audio, model, bahasa)};
  \node[kotakemas, minimum width=7.3cm, minimum height=24mm] (s) at (3.9,-1.7) {Panel status, log, dan hasil\\(tab Asli sampai Dokumen)};
  \node[kotak, minimum width=15cm, minimum height=10mm] (r) at (0,-3.6) {Panel riwayat (folder hasil milik akun Anda)};
  \draw[garis] (h) -- (u);
  \draw[garis] (h) -- (s);
\end{tikzpicture}
\caption{Tata letak kelompok panel pada layar utama Pandora}
\label{fig:tata-letak}
\end{figure}

Panel yang ditandai emas adalah tempat Anda membaca hasil, sedangkan panel unggah menjadi titik awal setiap pekerjaan.

\begin{table}[H]
\caption{Fungsi setiap panel}
\label{tab:panel}
{\small
\begin{tabularx}{\linewidth}{@{}P{4.2cm}Y@{}}
\toprule
\rowcolor{ugmpanel}\thd{Panel} & \thd{Fungsi} \\
\midrule
File audio & Memilih berkas rekaman dengan cara diseret atau diklik. \\
Pengaturan & Memilih model dan bahasa. Pilihan terakhir tersimpan. \\
Status dan log & Menampilkan progres, pesan, dan potongan teks secara langsung. \\
Hasil & Menampilkan tab Asli, JSON, Rapi AI, Notulensi, dan Dokumen. \\
Riwayat & Membuka kembali hasil lama dan menghapus hasil yang tidak diperlukan. \\
\bottomrule
\end{tabularx}
}
\end{table}

\begin{ilustrasi}
Rina membuka Pandora pada pagi hari dan melihat status layanan AI belum siap. Tombol untuk membuat notulensi berwarna pudar. Ia tidak menunda pekerjaan. Ia tetap menyiapkan rekaman pertama untuk ditranskripsi, lalu mencatat jam saat status itu ia lihat, supaya bisa ia laporkan kepada pengelola bila masih belum siap di siang hari.
\end{ilustrasi}

\begin{kegiatan}{Latihan 1. Masuk dan membaca layar (10 menit)}
Masuk dengan akun UGM Anda. Tulis status tiga komponen di bagian atas layar, lalu tentukan apakah tombol AI dapat dipakai hari ini. Jika tombol AI tidak aktif, catat penyebab yang paling mungkin menurut Tabel~\ref{tab:status}. Keluaran yang diharapkan: satu kalimat yang menyatakan kesiapan layanan dan satu kalimat tindakan yang akan Anda ambil.
\end{kegiatan}

\begin{rangkuman}
Pandora dibuka dengan akun UGM, dan hasil Anda hanya tampil di riwayat akun itu. Status layanan di bagian atas layar menentukan fitur mana yang siap dipakai, dan tombol AI tidak aktif bila layanan AI belum siap. Layar terdiri dari panel unggah, status dan hasil, serta riwayat. Pada bab berikutnya, Rina mulai dengan rekaman rapat pertamanya.
\end{rangkuman}
